\chapter{Lesson 5: Dining Out}

\section{Vocabulary}

\begin{center}
\begin{tabular}{llll}
\toprule
Characters & Pinyin & Part of speech & Meaning \\
\midrule
\zh{饭馆(儿)} & fànguǎn(r) & noun & restaurant \\
\zh{服务员} & fúwùyuán & noun & server, waiter \\
\zh{菜单} & càidān & noun & menu \\
\zh{点菜} & diǎn cài & verb-object & to order dishes \\
\zh{位} & wèi & measure word & (polite) for people \\
\zh{桌子} & zhuōzi & noun & table \\
\zh{饺子} & jiǎozi & noun & dumplings \\
\zh{米饭} & mǐfàn & noun & cooked rice \\
\zh{面条} & miàntiáo & noun & noodles \\
\zh{汤} & tāng & noun & soup \\
\zh{豆腐} & dòufu & noun & tofu \\
\zh{青菜} & qīngcài & noun & greens \\
\zh{牛肉} / \zh{猪肉} / \zh{鸡肉} & niúròu / zhūròu / jīròu & noun & beef / pork / chicken \\
\zh{鱼} & yú & noun & fish \\
\zh{素} & sù & adjective & vegetarian \\
\zh{酸} & suān & adjective & sour \\
\zh{甜} & tián & adjective & sweet \\
\zh{苦} & kǔ & adjective & bitter \\
\zh{辣} & là & adjective & spicy \\
\zh{咸} & xián & adjective & salty \\
\zh{饿} & è & adjective & hungry \\
\zh{渴} & kě & adjective & thirsty \\
\zh{饱} & bǎo & adjective & full \\
\zh{够} & gòu & adjective & enough \\
\zh{买单} & mǎi dān & verb-object & to pay the bill \\
\zh{打包} & dǎ bāo & verb-object & to pack leftovers \\
\zh{筷子} & kuàizi & noun & chopsticks \\
\zh{盘} & pán & measure word & plate of \\
\zh{碗} & wǎn & measure word & bowl of \\
\zh{又} & yòu & adverb & again (already happened) \\
\zh{再} & zài & adverb & again (not yet happened) \\
\zh{完} & wán & verb (complement) & to finish \\
\bottomrule
\end{tabular}
\end{center}

\section{In the restaurant}

\begin{itemize}
  \item \zh{请问，几位？}---\zh{三位。} How many people?---Three.
  \item \zh{你们想吃点儿什么？} What would you like to eat?
  \item \zh{来一盘饺子、一碗酸辣汤。} We'll have a plate of dumplings and a bowl of hot and
        sour soup. (\zh{来} = ``bring'' when ordering)
  \item \zh{请别放味精。} No MSG, please.
  \item \zh{我们吃素。} We are vegetarian.
  \item \zh{服务员，买单！} Server, the bill please!
  \item \zh{菜点多了，打包吧。} We ordered too much; let's pack it up.
\end{itemize}

\section{Grammar}

\subsection{Emphatic negation: \zh{一点儿也/都不}}

\zh{一点儿也不} + adjective = ``not \ldots at all''; \zh{一点儿也没} + verb = ``didn't \ldots
at all''.
\begin{itemize}
  \item \zh{这个菜一点儿也不辣。} This dish isn't spicy at all.
  \item \zh{我一点儿都不饿。} I'm not hungry at all.
  \item \zh{他一点儿饭也没吃。} He didn't eat a thing. (object between \zh{一点儿} and
        \zh{也})
\end{itemize}

\subsection{\zh{又} vs \zh{再}}

\begin{center}
\begin{tabular}{lp{5.8cm}p{5.8cm}}
\toprule
 & \zh{又} yòu & \zh{再} zài \\
\midrule
Time & repetition that already happened & repetition that will happen \\
Example & \zh{他昨天又去那个饭馆了。} & \zh{我们明天再去吧。} \\
 & He went to that restaurant again yesterday. & Let's go again tomorrow. \\
Requests & --- & \zh{请再来一碗米饭。} Another bowl of rice, please. \\
\bottomrule
\end{tabular}
\end{center}

\begin{remark}
\zh{又} is also used for predictable future repetitions: \zh{明天又是星期一了。}
``Tomorrow is Monday again.''
\end{remark}

\subsection{\zh{了} with a quantity}

When the object has a number, \zh{了} after the verb is natural and needed.
\begin{itemize}
  \item \zh{我吃了三碗饭。} I ate three bowls of rice.
  \item \zh{我们点了五个菜，太多了。} We ordered five dishes; too many.
\end{itemize}

\subsection{Resultative complements (introduction)}

A \emph{resultative complement} is a verb or adjective that follows the main verb and
says what \emph{result} the action had.

\begin{center}
\begin{tikzpicture}[w/.style={draw, rounded corners, minimum height=0.85cm, minimum width=1.4cm, align=center, font=\small}]
  \node[w, fill=orange!20] (v) at (0,0) {\zh{吃}\\ eat};
  \node[w, fill=green!20] (r) at (1.8,0) {\zh{完}\\ finish};
  \node[font=\small] at (3.3,0) {$\Rightarrow$};
  \node[font=\small, anchor=west] at (3.7,0) {\zh{吃完} ``eat up, finish eating''};
  \node[w, fill=orange!20] at (0,-1.2) {\zh{做}\\ do};
  \node[w, fill=green!20] at (1.8,-1.2) {\zh{好}\\ good};
  \node[font=\small] at (3.3,-1.2) {$\Rightarrow$};
  \node[font=\small, anchor=west] at (3.7,-1.2) {\zh{做好} ``do (it) properly, get it done''};
  \node[w, fill=orange!20] at (0,-2.4) {\zh{听}\\ listen};
  \node[w, fill=green!20] at (1.8,-2.4) {\zh{懂}\\ understand};
  \node[font=\small] at (3.3,-2.4) {$\Rightarrow$};
  \node[font=\small, anchor=west] at (3.7,-2.4) {\zh{听懂} ``understand by listening''};
  \node[w, fill=orange!20] at (0,-3.6) {\zh{写}\\ write};
  \node[w, fill=green!20] at (1.8,-3.6) {\zh{错}\\ wrong};
  \node[font=\small] at (3.3,-3.6) {$\Rightarrow$};
  \node[font=\small, anchor=west] at (3.7,-3.6) {\zh{写错} ``write incorrectly''};
  \node[font=\scriptsize] at (0,0.75) {action};
  \node[font=\scriptsize] at (1.8,0.75) {result};
\end{tikzpicture}
\end{center}

\begin{important}[Using resultative complements]
\begin{itemize}
  \item The object follows the complement: \zh{我吃完饭了}, not *\zh{我吃饭完了}.
  \item Negation uses \zh{没}: \zh{我还没吃完。} I haven't finished eating yet.
  \item Very common complements: \zh{完} (finish), \zh{好} (done properly), \zh{懂}
        (understand), \zh{见} (perceive: \zh{看见, 听见}), \zh{到} (reach the goal:
        \zh{找到, 买到}), \zh{错} (wrong), \zh{对} (correct), \zh{饱} (full: \zh{吃饱}).
\end{itemize}
\end{important}

\begin{itemize}
  \item \zh{你吃饱了吗？}---\zh{吃饱了，谢谢。} Are you full?---Yes, thanks.
  \item \zh{服务员上错菜了。} The server brought the wrong dish.
  \item \zh{我找到那家饭馆了。} I found that restaurant.
\end{itemize}

\section{Flavours and regional cuisines}

A well-known saying describes regional taste: \zh{南甜北咸，东辣西酸}
``sweet in the south, salty in the north, spicy in the east, sour in the west''. Modern
reality is messier (Sichuan in the west is famously spicy), so the class drew the map
as the main cuisines actually are.

\begin{center}
\begin{tikzpicture}[r/.style={draw, rounded corners, align=center, font=\small, fill=#1, minimum width=2.4cm, minimum height=1.1cm}]
  \draw[thick, gray, rounded corners=12pt] (-4.8,-2.6) -- (-5.4,0.6) -- (-3.6,3.3) -- (1,3.9) -- (4.9,3.3) -- (4.6,0.2) -- (3.4,-3.0) -- (0,-3.8) -- cycle;
  \node[r=blue!12] at (0.6,2.2) {\zh{北京/山东}\\ salty, wheat\\ \zh{饺子, 烤鸭}};
  \node[r=red!18] at (-2.8,-0.4) {\zh{四川/湖南}\\ \zh{麻辣} numbing-spicy\\ \zh{麻婆豆腐}};
  \node[r=yellow!25] at (2.6,-0.3) {\zh{上海/江苏}\\ sweet, delicate\\ \zh{小笼包}};
  \node[r=green!15] at (0.4,-2.4) {\zh{广东}\\ fresh, light\\ \zh{点心} dim sum};
  \node[r=orange!20] at (-2.9,2.0) {\zh{山西}\\ sour (vinegar)\\ \zh{刀削面}};
  \draw[->] (6.0,2.2) -- (6.0,3.0) node[above, font=\scriptsize] {\zh{北}};
\end{tikzpicture}
\end{center}

\begin{aside}
Table manners: never stick chopsticks upright in rice (it resembles incense at a
funeral); dishes are shared from the centre of the table; the host orders and keeps
serving guests. Fighting over the bill (\zh{抢着买单}) is a sign of friendship. Tipping
is not customary in mainland China, unlike Canada. Traditional characters often seen on
menus: \zht{餃子} \zh{饺子}, \zht{麵} \zh{面} (noodles), \zht{雞} \zh{鸡}.
\end{aside}

\section{Model conversation}

\begin{itemize}
  \item \textbf{Server:} \zh{几位？请坐。这是菜单。}
  \item \textbf{A:} \zh{我们来一盘饺子、一个麻婆豆腐和两碗米饭。}
  \item \textbf{B:} \zh{麻婆豆腐辣不辣？我一点儿辣的也不能吃。}
  \item \textbf{Server:} \zh{有一点儿辣。我们可以做得不辣。}
  \item \textbf{A:} \zh{好。再来一个青菜，一个酸辣汤。}
  \item \textbf{B:} \zh{点够了，别再点了！}
  \item \textbf{A:} \zh{（later）我吃饱了。菜没吃完，打包吧。服务员，买单！}
\end{itemize}

\section*{Review}
\begin{itemize}
  \item Ordering: \zh{来} + quantity + dish; \zh{几位？}; \zh{买单}.
  \item \zh{一点儿也/都不} + adj = not at all.
  \item \zh{又} for past repetition, \zh{再} for future repetition and requests.
  \item V + result: \zh{吃完, 做好, 听懂, 看见, 找到, 写错}; object after the
        complement; negate with \zh{没}.
  \item Flavours \zh{酸甜苦辣咸}.
\end{itemize}
